% ===========================================================================
% Paper 50 (DE): Algorithmische Zahlentheorie
% Specialist Mathematics Project
% @author: Michael Fuhrmann
% @date: 2026-03-12
% @build: 137
% ===========================================================================

\documentclass[12pt,a4paper]{amsart}
\usepackage{amsmath,amssymb,amsthm}
\usepackage{mathtools}
\usepackage[utf8]{inputenc}
\usepackage[T1]{fontenc}
\usepackage[ngerman]{babel}
\usepackage{hyperref}
\usepackage{enumitem,booktabs,array}
\usepackage{geometry}
\geometry{margin=2.5cm}

\newtheorem{theorem}{Satz}[section]
\newtheorem{lemma}[theorem]{Lemma}
\newtheorem{korollar}[theorem]{Korollar}
\newtheorem{proposition}[theorem]{Proposition}
\newtheorem{vermutung}[theorem]{Vermutung}
\theoremstyle{definition}
\newtheorem{definition}[theorem]{Definition}
\newtheorem{beispiel}[theorem]{Beispiel}
\newtheorem{bemerkung}[theorem]{Bemerkung}

\author{Michael Fuhrmann}
\address{Specialist Mathematics Project, Build 137}
\email{specialist-maths@localhost}
\date{\today}

\title{Algorithmische Zahlentheorie:\\
Grundlagen, Komplexit\"at und kryptographische Anwendungen}

\begin{document}

\maketitle
\tableofcontents

% ===========================================================================
\section{Einleitung}
% ===========================================================================

Die Zahlentheorie gehört zu den ältesten Gebieten der Mathematik, hat aber
durch die algorithmische Perspektive eine zentrale Rolle in der modernen
Informatik gewonnen. Die \emph{algorithmische Zahlentheorie} untersucht
ganzzahlige Probleme unter dem Gesichtspunkt ihrer Berechenbarkeit: Wie schnell
lässt sich prüfen, ob eine Zahl prim ist? Wie aufwändig ist die Faktorisierung
einer großen ganzen Zahl? Können diskrete Logarithmen effizient berechnet werden?

Diese Fragen sind nicht nur theoretischer Natur. Sie bilden das mathematische
Fundament der modernen Kryptographie. RSA-Verschlüsselung, Diffie-Hellman-
Schlüsselaustausch und elliptische-Kurven-Kryptosysteme schöpfen ihre Sicherheit
aus der angenommenen rechnerischen Schwierigkeit zahlentheoretischer Probleme.

\subsection*{Gliederung}

Wir behandeln fünf Kernthemen:
\begin{itemize}[leftmargin=2em]
  \item Größter gemeinsamer Teiler und erweiterter euklidischer Algorithmus
  \item Primzahltests: von probabilistisch bis deterministisch polynomiell
  \item Ganzzahlfaktorisierung und ihre Verbindung zur Kryptographie
  \item Algorithmen zum diskreten Logarithmus
  \item Chinesischer Restsatz, quadratische Reste, Kettenbrüche
\end{itemize}
Beweistheoretisch unterscheiden wir stets zwischen einem \emph{bewiesenen Satz}
(mit vollständigem Beweis) und einer \emph{Vermutung} (offen, kein Beweis bekannt).

\subsection*{Notation}

$\log$ bezeichnet den natürlichen Logarithmus. $\mathbb{Z}_n = \mathbb{Z}/n\mathbb{Z}$
ist der Restklassenring modulo $n$. $\varphi(n)$ bezeichnet Eulers Totient-Funktion.
$L_n[\alpha, c] = \exp\bigl(c(\log n)^\alpha(\log\log n)^{1-\alpha}\bigr)$ ist
die subexponentielle Komplexitätsnotation.

% ===========================================================================
\section{Größter gemeinsamer Teiler und erweiterter euklidischer Algorithmus}
% ===========================================================================

\subsection{Der euklidische Algorithmus}

Der \emph{größte gemeinsame Teiler} (GGT) zweier nicht-negativer ganzer Zahlen
$a$ und $b$, geschrieben $\gcd(a,b)$, ist die größte ganze Zahl, die beide teilt.
Der euklidische Algorithmus berechnet $\gcd(a,b)$ durch wiederholte Division:
\begin{align*}
  a &= q_0 b + r_0, \quad 0 \le r_0 < b,\\
  b &= q_1 r_0 + r_1, \quad 0 \le r_1 < r_0,\\
  &\vdots\\
  r_{k-2} &= q_k r_{k-1} + 0.
\end{align*}
Dann ist $\gcd(a,b) = r_{k-1}$, der letzte von Null verschiedene Rest.

\begin{theorem}[Euklidischer Algorithmus, Laufzeit. \textbf{Bewiesen.}]
Der euklidische Algorithmus terminiert nach $O(\log\min(a,b))$ Divisionsschritten.
Jeder Schritt halbiert den kleineren der beiden Reste (bis auf einen konstanten
Faktor), was eine geometrische Konvergenz garantiert.
\end{theorem}

\begin{proof}
Es gilt $r_{i+2} < r_i/2$: Falls $r_{i+1} \le r_i/2$, gilt dies direkt;
andernfalls ist $r_{i+2} = r_i - r_{i+1} < r_i/2$. Also halbiert sich der
Rest alle zwei Schritte, und nach $2\lfloor\log_2 \min(a,b)\rfloor$ Schritten
ist der Algorithmus beendet.
\end{proof}

\subsection{Erweiterter euklidischer Algorithmus und Bézout-Identität}

\begin{theorem}[Bézout-Identität. \textbf{Bewiesen.}]
Für ganze Zahlen $a,b$ (nicht beide null) existieren ganze Zahlen $s,t$ mit
\[
  sa + tb = \gcd(a,b).
\]
\end{theorem}

Der \emph{erweiterte euklidische Algorithmus} berechnet $\gcd(a,b)$ und die
Bézout-Koeffizienten $s,t$ gleichzeitig, indem er die Divisionsfolge
rückwärts verfolgt. Er läuft ebenfalls in $O(\log\min(a,b))$.

\subsubsection*{Modulares Inverses}

Das modulare Inverse von $a$ modulo $m$ (d.\,h. $x$ mit $ax \equiv 1 \pmod{m}$)
existiert genau dann, wenn $\gcd(a,m) = 1$, und ist gleich dem Bézout-Koeffizienten
$s$ reduziert modulo $m$. Dies wird durchgehend in RSA-Entschlüsselung und
CRT-Rekonstruktion verwendet.

\begin{beispiel}
Berechne $7^{-1} \pmod{11}$: Erweiterter Euklid liefert $1 = 8\cdot 7 - 5\cdot 11$,
also $7^{-1} \equiv 8 \pmod{11}$. Probe: $7 \cdot 8 = 56 = 5\cdot 11 + 1 \equiv 1$.
\end{beispiel}

\subsection{Lehrers Algorithmus und binärer GGT}

Für Multipräzisions-Arithmetik (Zahlen mit $k$ Maschinenwörtern) kostet
Schuldivision $O(k^2)$ Operationen. \emph{Lehrers Algorithmus} reduziert dies
auf $O(k\log^2 k \log\log k)$, indem er vollständige Divisionen durch
Einwort-Approximationen ersetzt. Der \emph{binäre GGT} vermeidet Divisionen
ganz durch Verschiebungen und Subtraktionen.

% ===========================================================================
\section{Primzahltests}
% ===========================================================================

\subsection{Probedivision}

Der einfachste Primzahltest: prüfe alle Zahlen von $2$ bis $\lfloor\sqrt{n}\rfloor$.
Teilt keine davon $n$, so ist $n$ prim. Laufzeit $O(\sqrt{n})$, exponentiell in
der Stellenzahl. Praktikabel nur für $n \lesssim 10^{12}$.

\subsection{Fermatscher Test und Pseudoprimzahlen}

\begin{theorem}[Kleiner Satz von Fermat. \textbf{Bewiesen.}]
Wenn $p$ eine Primzahl ist und $\gcd(a,p)=1$, so gilt $a^{p-1} \equiv 1 \pmod{p}$.
\end{theorem}

Prüfen von $a^{n-1} \equiv 1 \pmod{n}$ für mehrere Basen $a$ liefert eine
schnelle Primalitätsheuristik. Jedoch sind \emph{Carmichael-Zahlen} (z.\,B.
$561 = 3\cdot 11\cdot 17$) zusammengesetzte Zahlen, die $a^{n-1}\equiv 1\pmod{n}$
für alle $\gcd(a,n)=1$ erfüllen. Es gibt unendlich viele Carmichael-Zahlen
(bewiesen von Alford, Granville und Pomerance, 1994).

\subsection{Miller-Rabin-Test}

\begin{definition}[Starke Pseudoprimzahl]
Schreibe $n-1 = 2^s \cdot d$ mit $d$ ungerade. Dann ist $n$ eine
\emph{starke Pseudoprimzahl zur Basis $a$}, falls entweder
$a^d \equiv 1 \pmod{n}$ oder $a^{2^r d} \equiv -1 \pmod{n}$ für ein
$0 \le r < s$ gilt.
\end{definition}

\begin{theorem}[Miller--Rabin, 1976. \textbf{Bewiesen.}]
Wenn $n$ zusammengesetzt ist, ist der Anteil der Basen $a \in \{1,\ldots,n-1\}$,
für die $n$ eine starke Pseudoprimzahl zur Basis $a$ ist, höchstens $1/4$.
Nach $k$ unabhängigen zufälligen Basen fällt die Fehlerwahrscheinlichkeit auf
höchstens $4^{-k}$.
\end{theorem}

Mit $k = 40$ Runden sinkt die Fehlerwahrscheinlichkeit unter $10^{-24}$, was
für alle praktischen Zwecke ausreicht. Die deterministische Variante mit festen
Basen ist für alle $n < 3{,}3 \times 10^{24}$ beweisbar korrekt.

\subsection{Solovay-Strassen-Test}

\begin{theorem}[Eulerkriterium. \textbf{Bewiesen.}]
Für eine ungerade Primzahl $p$ und $\gcd(a,p)=1$:
$a^{(p-1)/2} \equiv \bigl(\frac{a}{p}\bigr) \pmod{p}$,
wobei $\bigl(\frac{a}{p}\bigr)$ das Legendre-Symbol ist.
\end{theorem}

Der \emph{Solovay-Strassen-Test} verallgemeinert dies über das Jacobi-Symbol:
Prüfe $a^{(n-1)/2} \equiv \bigl(\frac{a}{n}\bigr)_J \pmod{n}$. Es gibt keine
analogen universellen Gegenbeispiele wie Carmichael-Zahlen beim Fermat-Test:
Für zusammengesetztes $n$ besteht höchstens die Hälfte aller Basen den Test.

\subsection{AKS-Primzahltest}

\begin{theorem}[Agrawal--Kayal--Saxena, 2002. \textbf{Bewiesen.}]
\textit{PRIMES liegt in P.} Es gibt einen deterministischen Algorithmus, der
die Primalität von $n$ in Zeit $\tilde{O}(\log^6 n)$ entscheidet.
\end{theorem}

Der AKS-Test beruht auf der Polynomidentität: $n$ ist genau dann prim, wenn
\[
  (X+a)^n \equiv X^n + a \pmod{X^r - 1,\; n}
\]
für alle $a \le \lfloor\sqrt{\varphi(r)}\log n\rfloor$ gilt, wobei $r$ so
gewählt wird, dass die multiplikative Ordnung von $n$ modulo $r$ den Wert
$\log^2 n$ übersteigt.

\subsubsection*{Algorithmusschritte}
\begin{enumerate}[label=\arabic*.]
  \item \textbf{Perfekte-Potenz-Test:} Falls $n = a^b$ mit $b \ge 2$,
    gib \textsc{Zusammengesetzt} aus.
  \item \textbf{$r$ bestimmen:} Finde das kleinste $r$ mit
    $\operatorname{ord}_r(n) > \log^2 n$.
  \item \textbf{Kleiner-Teiler-Test:} Falls $1 < \gcd(a,n) < n$ für ein
    $a \le r$, gib \textsc{Zusammengesetzt} aus.
  \item \textbf{Größenprüfung:} Falls $n \le r$, gib \textsc{Prim} aus.
  \item \textbf{Polynomtest:} Prüfe $(X+a)^n \equiv X^n + a \pmod{X^r-1, n}$
    für $a = 1, \ldots, \lfloor\sqrt{\varphi(r)}\log n\rfloor$.
  \item Gib \textsc{Prim} aus.
\end{enumerate}

\begin{bemerkung}
Die ursprüngliche Komplexität $O(\log^{12} n)$ wurde von Lenstra und Pomerance
(2005) auf $O(\log^6 n)$ verbessert. AKS beweist, dass \textsc{PRIMES} $\in$ P,
obwohl Miller-Rabin in der Praxis schneller bleibt.
\end{bemerkung}

% ===========================================================================
\section{Ganzzahlfaktorisierung}
% ===========================================================================

\subsection{Das Faktorisierungsproblem}

\begin{vermutung}[Schwierigkeit der Faktorisierung. \textbf{Offen.}]
Kein Polynomialzeit-Algorithmus für die ganzzahlige Faktorisierung (IFP) ist bekannt.
Insbesondere ist IFP vermutlich nicht in P. Die Verbindung zu P $\ne$ NP ist subtil:
IFP liegt in NP $\cap$ co-NP, wäre also nicht NP-vollständig, es sei denn
NP $=$ co-NP. Der beste bekannte klassische Algorithmus (GNFS) läuft in
subexponentieller Zeit $L_n[1/3, c]$.
\end{vermutung}

Diese Vermutung ist das Fundament der RSA-Sicherheit. Kein polynomieller
Algorithmus ist bekannt.

\subsection{Pollards $\rho$-Algorithmus}

\emph{Pollards $\rho$-Algorithmus} (1975) findet einen nicht-trivialen Teiler
von $n$ über eine pseudozufällige Folge $x_{i+1} = f(x_i) \bmod n$
(typischerweise $f(x) = x^2 + c$). Durch das Geburtstagsproblem ist eine
Kollision $\gcd(|x_i - x_j|, n) > 1$ nach $O(n^{1/4})$ Schritten zu erwarten.
Floyds Zykluserkennung reduziert den Speicherbedarf auf $O(1)$.

\begin{theorem}[Erwartete Komplexität von Pollard $\rho$. Heuristisch.]
Unter heuristischen Annahmen zur Pseudozufälligkeit von $f$ ist die erwartete
Schrittzahl zum Finden eines Teilers $p \mid n$ gleich $O(p^{1/2}) = O(n^{1/4})$.
\end{theorem}

\subsection{Fermat-Faktorisierung und Quadratisches Sieb}

\emph{Fermats Methode} nutzt $n = a^2 - b^2 = (a-b)(a+b)$: Suche
$a \ge \lceil\sqrt{n}\rceil$ mit $a^2 - n$ als perfekte Quadratzahl. Effizient
wenn $n$ zwei annähernd gleich große Primteiler hat.

Das \emph{Quadratische Sieb} (Pomerance, 1981) ist für $n \lesssim 10^{100}$
der schnellste bekannte Algorithmus. Es sucht Paare $(x,y)$ mit
$x^2 \equiv y^2 \pmod{n}$, dann liefert $\gcd(x-y, n)$ einen Teiler.
Die Siebphase findet viele $B$-glatte Werte von $x^2 \bmod n$ gleichzeitig
und erreicht heuristisch $L_n[1/2, 1]$.

\subsection{Das allgemeine Zahlkörpersieb (GNFS)}

Das \emph{General Number Field Sieve} (GNFS) ist der aktuell schnellste bekannte
Algorithmus für $n > 10^{100}$. Es operiert in einem algebraischen Zahlkörper
$\mathbb{Z}[\theta]$ für eine Wurzel $\theta$ eines sorgsam gewählten Polynoms
und siebt nach glatten Normen.

\begin{theorem}[GNFS-Komplexität. Heuristisch.]
GNFS faktorisiert $n$ in erwarteter Zeit
\[
  L_n\!\left[\tfrac{1}{3},\, \left(\tfrac{64}{9}\right)^{1/3}\right]
  \approx \exp\!\left(1{,}923\,(\log n)^{1/3}(\log\log n)^{2/3}\right).
\]
\end{theorem}

Der 512-Bit-RSA-155-Rekord (1999) benötigte 8000 MIPS-Jahre; der 829-Bit-RSA-250-Rekord
(2020) erforderte etwa 2700 CPU-Kernjahre.

\subsection{Implikationen für RSA}

\begin{vermutung}[RSA-Sicherheit. \textbf{Offen.}]
Gegeben ein Produkt $n = pq$ zweier großer Primzahlen und ein Chiffretext
$c = m^e \bmod n$: Die Rückgewinnung von $m$ ohne Kenntnis von $p,q$ gilt als
rechnerisch nicht machbar. Kein Polynomialzeit-Algorithmus für das RSA-Problem
ist bekannt. Das RSA-Problem ist polynomial-Zeit-äquivalent zur Faktorisierung
von $n$ unter natürlichen Reduktionen.
\end{vermutung}

Aktuelle Empfehlung: RSA-Schlüssellängen $\ge 3072$ Bit (NIST 2022) oder
Umstieg auf elliptische-Kurven-Systeme.

% ===========================================================================
\section{Das Problem des diskreten Logarithmus}
% ===========================================================================

\subsection{Definition und Schwierigkeit}

\begin{definition}[Diskreter Logarithmus]
Sei $g$ ein Erzeuger der zyklischen Gruppe $(\mathbb{Z}/p\mathbb{Z})^*$ der
Ordnung $p-1$. Der \emph{diskrete Logarithmus} von $h$ zur Basis $g$ ist das
eindeutige $x \in \{0,1,\ldots,p-2\}$ mit $g^x \equiv h \pmod{p}$.
\end{definition}

\begin{vermutung}[Schwierigkeit des DLP. \textbf{Offen.}]
Kein Polynomialzeit-Algorithmus für das diskrete-Logarithmus-Problem (DLP) in
$(\mathbb{Z}/p\mathbb{Z})^*$ ist bekannt. Die besten Algorithmen erreichen
subexponentielle Komplexität $L_p[1/3, c]$, dieselbe wie GNFS.
\end{vermutung}

\subsection{Baby-Step-Giant-Step (BSGS)}

\begin{theorem}[Shanks, 1971. \textbf{Bewiesen.}]
Der Baby-Step-Giant-Step-Algorithmus berechnet $\log_g h$ in
$(\mathbb{Z}/p\mathbb{Z})^*$ in $O(\sqrt{p})$ Zeit und Speicher.
\end{theorem}

Setze $m = \lceil\sqrt{p-1}\rceil$ und schreibe $x = im - j$:
\begin{itemize}
  \item \textbf{Baby-Schritte:} Berechne die Tabelle
    $\{(h \cdot g^j \bmod p, j) : 0 \le j < m\}$.
  \item \textbf{Giant-Schritte:} Berechne $g^{im} \bmod p$ für $i = 1,\ldots,m$.
  \item \textbf{Kollision:} Wenn $g^{im} = h \cdot g^j$, gilt $x = im - j$.
\end{itemize}

\subsection{Pohlig-Hellman-Algorithmus}

\begin{theorem}[Pohlig--Hellman, 1978. \textbf{Bewiesen.}]
Wenn $p-1 = \prod_i q_i^{e_i}$ mit allen $q_i$ klein ist, kann der diskrete
Logarithmus in $(\mathbb{Z}/p\mathbb{Z})^*$ in
$O\bigl(\sum_i e_i(\log(p-1) + \sqrt{q_i})\bigr)$ Operationen berechnet werden.
\end{theorem}

Der Algorithmus löst das DLP modulo jeder Primpotenz $q_i^{e_i}$ separat mittels
BSGS in der entsprechenden Untergruppe, kombiniert dann via CRT. Dies macht das
DLP leicht, wenn $p-1$ \emph{glatt} ist (alle Primfaktoren klein).

\begin{bemerkung}
Dies motiviert die Verwendung \emph{sicherer Primzahlen} $p = 2q+1$ (mit $q$
ebenfalls prim) in Diffie-Hellman: Dann hat $p-1 = 2q$ den großen Primteiler $q$,
was Pohlig-Hellman wirkungslos macht.
\end{bemerkung}

\subsection{Indexkalkül}

Der \emph{Indexkalkül-Algorithmus} erreicht subexponentielle Komplexität für
das DLP in $\mathbb{F}_p^*$. Er läuft in zwei Phasen:

\begin{enumerate}[label=\arabic*.]
  \item \textbf{Beziehungsphase:} Aufbau einer Faktorbasis
    $\mathcal{B} = \{p_1,\ldots,p_k\}$ kleiner Primzahlen. Für zufällige
    Exponenten $s$: Prüfe ob $g^s \bmod p$ glatt ist. Falls ja, notiere
    $s \equiv \sum_i a_i \log_g p_i \pmod{p-1}$.
  \item \textbf{Individueller Logarithmus:} Finde $r$ mit $h \cdot g^r$ glatt,
    berechne $\log_g h$.
\end{enumerate}

\begin{theorem}[Indexkalkül-Komplexität. Heuristisch.]
Indexkalkül berechnet diskrete Logarithmen in $(\mathbb{Z}/p\mathbb{Z})^*$
heuristisch in Zeit $L_p[1/2, 1]$; mit Zahlkörpersieb-Variante in $L_p[1/3, c]$.
\end{theorem}

% ===========================================================================
\section{Chinesischer Restsatz}
% ===========================================================================

\subsection{Aussage und Beweis}

\begin{theorem}[Chinesischer Restsatz (CRS). \textbf{Bewiesen.}]
Seien $m_1, m_2, \ldots, m_k$ paarweise teilerfremde positive ganze Zahlen mit
$M = m_1 m_2 \cdots m_k$. Für beliebige ganze Zahlen $r_1, r_2, \ldots, r_k$
hat das System
\[
  x \equiv r_i \pmod{m_i}, \quad i = 1, \ldots, k,
\]
eine eindeutige Lösung $x \pmod{M}$.
\end{theorem}

\begin{proof}
Setze $M_i = M/m_i$ und $e_i = M_i \cdot (M_i^{-1} \bmod m_i)$. Dann gilt
$e_i \equiv 1 \pmod{m_i}$ und $e_i \equiv 0 \pmod{m_j}$ für $j \ne i$.
Also erfüllt $x = \sum_{i=1}^k r_i e_i \bmod M$ alle Kongruenzen. Die
Eindeutigkeit folgt daraus, dass zwei Lösungen sich um ein Vielfaches von $M$
unterscheiden.
\end{proof}

\subsection{Algorithmische Komplexität}

Die CRS-Rekonstruktion läuft in $O(k \log^2 M)$ Bitoperationen (mit schneller
modularer Arithmetik). Für konstantes $k$ ergibt sich $O(\log^2 M)$, d.\,h.
quasi-linear in der Stellenzahl.

\subsection{Kryptographische Anwendungen}

\subsubsection*{RSA-CRT-Entschlüsselung}

Statt $c^d \bmod n$ direkt zu berechnen ($d \approx n$), berechnet man
$c^{d_p} \bmod p$ und $c^{d_q} \bmod q$ separat
(mit $d_p = d \bmod (p-1)$, $d_q = d \bmod (q-1)$) und kombiniert via CRS.
Dies liefert einen etwa $4$-fachen Geschwindigkeitsvorteil (Garners Formel).

\subsubsection*{Multimodulare Arithmetik}

Große Ganzzahlmultiplikation kann durch Arbeiten modulo mehrerer kleiner
Primzahlen und anschließender CRS-Rekonstruktion ausgeführt werden, was
Multipräzisions-Arithmetik vermeidet.

% ===========================================================================
\section{Quadratische Reste und der Tonelli-Shanks-Algorithmus}
% ===========================================================================

\subsection{Legendre-Symbol und Eulerkriterium}

\begin{definition}[Legendre-Symbol]
Für eine ungerade Primzahl $p$ und $\gcd(a,p)=1$:
\[
  \left(\frac{a}{p}\right) =
  \begin{cases} 1 & \text{falls } a \text{ ein quadratischer Rest mod } p \text{ ist},\\
                -1 & \text{sonst.}
  \end{cases}
\]
\end{definition}

\begin{theorem}[Eulerkriterium. \textbf{Bewiesen.}]
$\left(\frac{a}{p}\right) \equiv a^{(p-1)/2} \pmod{p}$.
\end{theorem}

\subsection{Gaußsches quadratisches Reziprozitätsgesetz}

\begin{theorem}[Gauß, 1796. \textbf{Bewiesen.}]
Für verschiedene ungerade Primzahlen $p, q$:
\[
  \left(\frac{p}{q}\right)\left(\frac{q}{p}\right) =
  (-1)^{\tfrac{p-1}{2}\cdot\tfrac{q-1}{2}}.
\]
Äquivalent: $\bigl(\frac{p}{q}\bigr) = \bigl(\frac{q}{p}\bigr)$,
außer wenn $p \equiv q \equiv 3 \pmod{4}$; dann haben sie entgegengesetzte Vorzeichen.
\end{theorem}

Gauß gab acht Beweise. Zusammen mit dem ersten und zweiten Ergänzungssatz
(für $(-1/p)$ und $(2/p)$) erlaubt das Reziprozitätsgesetz die effiziente
Berechnung des Legendre-Symbols in $O(\log^2 p)$ Schritten.

\subsection{Jacobi-Symbol}

Das \emph{Jacobi-Symbol} $\bigl(\frac{a}{n}\bigr)_J$ erweitert das Legendre-Symbol
auf zusammengesetzte $n$ via Multiplikativität über die Primfaktorzerlegung von $n$.
Es ist in $O(\log^2 n)$ ohne Kenntnis der Faktorisierung von $n$ berechenbar.
Jedoch gilt: $\bigl(\frac{a}{n}\bigr)_J = 1$ impliziert \emph{nicht}, dass $a$
ein QR modulo $n$ ist, wenn $n$ zusammengesetzt ist.

\subsection{Tonelli-Shanks-Algorithmus}

Gegeben einem quadratischen Rest $a$ modulo einer ungeraden Primzahl $p$,
berechnet der \emph{Tonelli-Shanks-Algorithmus} $\sqrt{a} \bmod p$ (d.\,h.
$x$ mit $x^2 \equiv a \pmod p$) in erwarteter $O(\log^2 p)$-Zeit.

Schreibe $p-1 = 2^s \cdot q$ mit $q$ ungerade. Finde ein Nicht-Rest $z$ und
initialisiere $M = s$, $c = z^q \bmod p$, $t = a^q \bmod p$, $R = a^{(q+1)/2} \bmod p$.
Der Algorithmus iteriert, halbiert $M$ und passt $R$ via $c$ an, bis $t = 1$.

\begin{theorem}[Tonelli-Shanks-Komplexität. \textbf{Bewiesen.}]
Die erwartete Laufzeit beträgt $O(\log^2 p)$ modulare Multiplikationen, wobei
die Erwartung über die Zufallswahl des Nicht-Restes $z$ genommen wird.
\end{theorem}

% ===========================================================================
\section{Kettenbrüche}
% ===========================================================================

\subsection{Definition und periodische Entwicklung}

\begin{definition}[Kettenbruch]
Ein \emph{einfacher Kettenbruch} ist ein Ausdruck
$[a_0; a_1, a_2, \ldots] = a_0 + \cfrac{1}{a_1 + \cfrac{1}{a_2 + \cdots}}$
mit $a_i \in \mathbb{Z}$, $a_i > 0$ für $i \ge 1$.
\end{definition}

\begin{theorem}[Lagrange, 1770. \textbf{Bewiesen.}]
Die Kettenbruchentwicklung von $\sqrt{n}$ (für $n$ keine Quadratzahl) ist
periodisch: $\sqrt{n} = [a_0; \overline{a_1, a_2, \ldots, a_k}]$ mit Periodenlänge
$k \le 2a_0$.
\end{theorem}

Die Konvergenten $p_k/q_k$ von $\sqrt{n}$ erfüllen
$|p_k - q_k\sqrt{n}| < 1/q_{k+1}$ — sie sind die besten rationalen
Approximationen an $\sqrt{n}$.

\subsection{Pell-Gleichung}

\begin{theorem}[Lagrange. \textbf{Bewiesen.}]
Die Gleichung $x^2 - ny^2 = 1$ (Pell-Gleichung, $n$ keine Quadratzahl) hat
unendlich viele ganzzahlige Lösungen. Die Fundamentallösung $(x_1, y_1)$
wird durch die Konvergenten von $\sqrt{n}$ geliefert.
\end{theorem}

\begin{beispiel}
$n=2$: $\sqrt{2} = [1;\overline{2}]$, Periode $k=1$. Fundamentallösung:
$(p_1, q_1) = (3, 2)$: $3^2 - 2\cdot 2^2 = 9-8 = 1$.
\end{beispiel}

\subsection{CFRAC-Faktorisierung}

Der \emph{Kettenbruch-Faktorisierungsalgorithmus} (Morrison--Brillhart, 1975)
ist ein frühes Glatt-Zahlen-Verfahren. Er sucht Paare $(x,y)$ mit
$x^2 \equiv y^2 \pmod{n}$, indem er die Zähler $p_k$ der Konvergenten
von $\sqrt{n}$ sammelt und auf $B$-Glattheit von $p_k^2 \bmod n$ siebt.
Dies war einer der ersten Algorithmen mit subexponentieller Heuristik-Komplexität.

\subsection{Kettenbrüche und Kryptoanalyse}

Kettenbrüche erscheinen im \emph{Wiener-Angriff auf RSA}: Falls der geheime
Exponent $d < n^{1/4}$, kann $d$ aus der Kettenbruchentwicklung von $e/n$
zurückgewonnen werden (Wiener, 1990).

\begin{bemerkung}
Der Rekursionsalgorithmus für Konvergenten $p_k/q_k$:
$p_k = a_k p_{k-1} + p_{k-2}$, $q_k = a_k q_{k-1} + q_{k-2}$,
läuft in $O(k)$ Zeit und ist numerisch sehr stabil.
\end{bemerkung}

% ===========================================================================
\section{Kryptographische Anwendungen}
% ===========================================================================

\subsection{RSA-Kryptosystem}

\begin{definition}[RSA, Rivest--Shamir--Adleman, 1978]
\textbf{Schlüsselerzeugung:} Wähle große Primzahlen $p,q$, setze $n = pq$,
$\varphi(n) = (p-1)(q-1)$. Wähle öffentlichen Exponenten $e$ mit
$\gcd(e, \varphi(n)) = 1$; berechne $d = e^{-1} \bmod \varphi(n)$ via
Erweiterter Euklid. \textbf{Öffentlicher Schlüssel:} $(n, e)$.
\textbf{Geheimer Schlüssel:} $d$.
\textbf{Verschlüsselung:} $c = m^e \bmod n$.
\textbf{Entschlüsselung:} $m = c^d \bmod n$ (folgt aus $ed \equiv 1 \pmod{\varphi(n)}$
und dem Satz von Euler).
\end{definition}

Die RSA-Sicherheit beruht auf der Schwierigkeit der Faktorisierung von $n$
(vgl.\ Abschnitt~4.4).

\subsection{Diffie-Hellman-Schlüsselaustausch}

\begin{definition}[DH-Protokoll, 1976]
Öffentlich: große Primzahl $p$, Erzeuger $g$ von $(\mathbb{Z}/p\mathbb{Z})^*$.
Alice sendet $g^a \bmod p$, Bob sendet $g^b \bmod p$. Beide berechnen das
gemeinsame Geheimnis $g^{ab} \bmod p$. Ein Angreifer muss das DLP lösen,
um $a$ oder $b$ zu erhalten.
\end{definition}

\begin{vermutung}[Computational Diffie-Hellman, CDH. \textbf{Offen.}]
Gegeben $g^a$ und $g^b$ in $(\mathbb{Z}/p\mathbb{Z})^*$: Die Berechnung von
$g^{ab}$ ist schwer. Die CDH-Annahme ist mindestens so schwer wie DLP.
\end{vermutung}

\subsection{Elliptische-Kurven-Kryptographie (ECC)}

Sei $E: y^2 = x^3 + ax + b$ über $\mathbb{F}_p$. Die Gruppe $E(\mathbb{F}_p)$
unter dem Sekanten-Tangenten-Gesetz hat Ordnung $\approx p$ (Satz von Hasse).
Das ECDLP (Elliptic Curve DLP) gilt als schwerer als das klassische DLP bei
gleicher Gruppengröße, da kein subexponentieller allgemeiner Algorithmus für
elliptische Kurven bekannt ist.

\begin{theorem}[Hasse, 1933. \textbf{Bewiesen.}]
$\bigl||E(\mathbb{F}_p)| - (p+1)\bigr| \le 2\sqrt{p}$.
\end{theorem}

ECDSA und ECDH verwenden 256-Bit-Schlüssel für eine Sicherheit äquivalent zu
RSA-3072, was einen erheblichen Effizienzgewinn darstellt.

\subsection{Post-Quanten-Kryptographie}

Shors Quantenalgorithmus (1994) löst sowohl Ganzzahlfaktorisierung als auch DLP
in $O((\log n)^3)$ Quantenzeit und bricht RSA, DH und ECC. Dies motiviert die
\emph{Post-Quanten-Kryptographie}:

\begin{itemize}
  \item \textbf{Gitterbasiert:} CRYSTALS-Kyber (NIST PQC Standard, 2022)
  \item \textbf{Hashbasiert:} SPHINCS+
  \item \textbf{Codebasiert:} McEliece (1978)
  \item \textbf{Multivariat:} HQC
\end{itemize}

NIST standardisierte 2022--2024 die ersten Post-Quanten-Algorithmen. Die
mathematische Schwierigkeit von Gitterproblemen (LWE, NTRU) beruht auf
vermuteten Worst-Case-auf-Average-Case-Reduktionen.

\subsection{Zusammenfassung der kryptographischen Annahmen}

\begin{table}[h]
\centering
\begin{tabular}{lll}
\toprule
\textbf{Kryptosystem} & \textbf{Mathematisches Problem} & \textbf{Status} \\
\midrule
RSA & Ganzzahlfaktorisierung (IFP) & Offen (Vermutung) \\
Diffie-Hellman & Diskreter Logarithmus (DLP) & Offen (Vermutung) \\
ECDH/ECDSA & Elliptischer-Kurven-DLP & Offen (Vermutung) \\
AKS & PRIMES $\in$ P & \textbf{Bewiesen} (2002) \\
CRS-Korrektheit & Zahlentheoretisches Theorem & \textbf{Bewiesen} \\
Reziprozitätsgesetz & Legendres/Jacobis Symbole & \textbf{Bewiesen} \\
Pell-Gleichung & Kettenbrüche und $\sqrt{n}$ & \textbf{Bewiesen} \\
\bottomrule
\end{tabular}
\caption{Kryptosysteme, ihre mathematischen Grundlagen und Beweisstatus}
\end{table}

% ===========================================================================
\section{Schlussbetrachtung und Verbindungen}
% ===========================================================================

Die algorithmische Zahlentheorie steht an der Kreuzung von reiner Mathematik
und theoretischer Informatik. Die Kernergebnisse gliedern sich in drei Bereiche:

\medskip
\noindent\textbf{Bewiesene Sätze} als Grundlage der Algorithmen:
Euklidischer Algorithmus ($O(\log n)$ Schritte), Bézout-Identität, Fermatscher
Kleiner Satz, AKS-Primalität in P, Gaußsches Reziprozitätsgesetz, Lagrange zur
Pell-Gleichung, Hassescher Satz zu elliptischen Kurven.

\medskip
\noindent\textbf{Allgemein geglaubte, aber unbewiesene Vermutungen}:
Schwierigkeit der Faktorisierung (IFP $\notin$ P), DLP-Schwierigkeit (CDH)
und RSA-Sicherheit. Dies sind die tragenden Annahmen der modernen Kryptographie.
Die Verbindung zu P $\ne$ NP ist indirekt: IFP $\in$ NP $\cap$ co-NP, wäre also
nicht NP-vollständig, es sei denn die polynomielle Hierarchie kollabiert.

\medskip
\noindent\textbf{Offene Probleme}:
Zwillingsprimzahlvermutung, Goldbachsche Vermutung, Verteilung von Primzahllücken
(Cramérs Vermutung), effiziente Algorithmen für ECDLP.

Der Übergang zur Post-Quanten-Kryptographie, getrieben durch Shors Algorithmus,
illustriert, wie Fortschritte in der algorithmischen Zahlentheorie die
Sicherheitslandschaft globaler Kommunikation direkt umgestalten.

% ===========================================================================
\begin{thebibliography}{99}
% ===========================================================================

\bibitem{AKS2002}
M.~Agrawal, N.~Kayal, N.~Saxena,
\emph{PRIMES is in P},
Ann.\ of Math.\ \textbf{160} (2004), 781--793.

\bibitem{CP2005}
R.~Crandall, C.~Pomerance,
\emph{Prime Numbers: A Computational Perspective},
2.~Aufl., Springer, New York, 2005.

\bibitem{Cohen1993}
H.~Cohen,
\emph{A Course in Computational Algebraic Number Theory},
Graduate Texts in Mathematics \textbf{138}, Springer, 1993.

\bibitem{MenezesOV1996}
A.~Menezes, P.~van Oorschot, S.~Vanstone,
\emph{Handbook of Applied Cryptography},
CRC Press, 1996. (\url{https://cacr.uwaterloo.ca/hac/}).

\bibitem{Shoup2009}
V.~Shoup,
\emph{A Computational Introduction to Number Theory and Algebra},
2.~Aufl., Cambridge University Press, 2009.

\bibitem{Pollard1975}
J.~M.~Pollard,
\emph{A Monte Carlo method for factorization},
BIT Numerical Mathematics \textbf{15} (1975), 331--334.

\bibitem{Miller1976}
G.~L.~Miller,
\emph{Riemann's hypothesis and tests for primality},
J.~Comput.\ System Sci.\ \textbf{13} (1976), 300--317.

\bibitem{PohligHellman1978}
S.~C.~Pohlig, M.~E.~Hellman,
\emph{An improved algorithm for computing logarithms over $GF(p)$ and its cryptographic significance},
IEEE Trans.\ Inform.\ Theory \textbf{24} (1978), 106--110.

\bibitem{Shanks1971}
D.~Shanks,
\emph{Class Number, a Theory of Factorization and Genera},
Proc.\ Symp.\ Pure Math.\ \textbf{20}, AMS, 1971, S.~415--440.

\bibitem{Shor1994}
P.~W.~Shor,
\emph{Algorithms for quantum computation: discrete logarithms and factoring},
Proc.\ 35th FOCS, 1994, S.~124--134.

\end{thebibliography}

\end{document}
